% Sections supplémentaires - inclus par REVISION_PSY0.tex
% ETOPS, CRM, Carburant, Stratégie AF, Systèmes avion, Balisage piste

%=============================================================
\section{ETOPS \& OPÉRATIONS ÉTENDUES}
%=============================================================

\subsection{Définition}
\textbf{ETOPS} = \textit{Extended-range Twin-engine Operational Performance Standards} (ou \textbf{EDTO} depuis 2012, OACI).\\
Règle qui permet aux \textbf{biréacteurs} d'opérer sur des routes éloignées d'un aéroport de déroutement. Le chiffre ETOPS = temps max \textbf{en minutes} avec \textbf{1 moteur en panne} pour rejoindre un aéroport.

\subsection{Niveaux de Certification}

\begin{tabular}{p{3cm}p{3cm}p{8.5cm}}
\toprule
\textbf{ETOPS} & \textbf{Couverture} & \textbf{Exemples} \\
\midrule
60 min & Routes courtes & Règle historique, avant ETOPS \\
120 min & Atlantique partiel & A300, premières certifications \\
\rowcolor{aflight}
\textbf{180 min (3h)} & \textbf{$\sim$95\% du globe} & A330, B767, B777, A320neo \\
240 min & 99\% du globe & B777-200LR \\
\rowcolor{aflight}
\textbf{370 min (6h10)} & \textbf{99.7\% du globe} & \textbf{A350-900} (1er certifié, 2014) \\
\bottomrule
\end{tabular}

\begin{memobox}
\textbf{Pourquoi l'ETOPS est important ?} Avant cette règle, les biréacteurs ne pouvaient pas traverser l'Atlantique~! Seuls les quadriréacteurs (A340, B747) le pouvaient. Grâce à l'ETOPS-180+, le B777 a rendu les quadriréacteurs obsolètes (plus chers, plus gourmands).
\end{memobox}

%=============================================================
\section{CARBURANT AVIATION}
%=============================================================

\subsection{Jet A-1 --- Fiche Technique}

\begin{tabular}{p{4.5cm}p{10cm}}
\toprule
\textbf{Type} & Kérosène aviation \\
Composition & Hydrocarbures C9--C16 (paraffines, aromatiques) \\
\textbf{Densité} & $\sim$\textbf{0.8 kg/L} (775--840 kg/m$^3$) \\
\textbf{Point d'éclair} & $>$ \textbf{38°C} (sécurité incendie) \\
\textbf{Point de congélation} & $\leq$ \textbf{$-$47°C} (vol en altitude) \\
Pouvoir calorifique & 43.15 MJ/kg \\
Ébullition & 150--250°C \\
Norme & DEF STAN 91-91, ASTM D1655 \\
\bottomrule
\end{tabular}

\subsection{SAF --- Carburant Aviation Durable}
\begin{itemize}[nosep]
\item SAF = \textit{Sustainable Aviation Fuel}, produit à partir de biomasse, huiles usagées, déchets
\item Réduit le CO$_2$ de 50--80\% sur le cycle de vie complet
\item Objectif AF~: \textbf{10\% de SAF} d'ici 2030
\item Mélangeable au Jet A-1 (max 50\% actuellement)
\item Pas de modification moteur nécessaire (``\textit{drop-in fuel}'')
\end{itemize}

\subsection{Avgas 100LL}
Carburant pour avions \textbf{à pistons} (avions légers). Essence contenant encore du plomb (coloration bleue). 100 octanes.

%=============================================================
\section{FACTEURS HUMAINS \& CRM}
%=============================================================

\subsection{CRM --- Crew Resource Management}

\begin{tabular}{p{4cm}p{10.5cm}}
\toprule
\textbf{Définition} & Formation aux compétences non-techniques pour améliorer la sécurité \\
Origine & Atelier NASA 1979, après crash Tenerife 1977 (583 morts) et UA 173 (1978) \\
Objectif & Gestion optimale de \textbf{toutes les ressources} disponibles (humaines, matérielles, informationnelles) \\
\bottomrule
\end{tabular}

\subsection{Compétences CRM}
\begin{enumerate}[nosep]
\item \textbf{Communication} --- phraséologie standard, read-back/hear-back
\item \textbf{Conscience situationnelle} --- modèle mental partagé de la situation
\item \textbf{Prise de décision} --- gestion du temps, évaluation des options
\item \textbf{Gestion de la charge de travail} --- priorisation, délégation
\item \textbf{Leadership et assertivité} --- le copilote \textbf{doit} dire s'il voit un problème
\item \textbf{Travail d'équipe} --- synergie PNT + PNC + ATC
\end{enumerate}

\subsection{Modèle de Reason --- ``Fromage Suisse''}

\begin{center}
\fbox{\parbox{13cm}{
Chaque \textbf{couche de défense} (procédures, formation, technologie, réglementation) est comme une tranche de fromage suisse avec des \textbf{trous} (faiblesses). Un accident survient quand les trous de \textbf{toutes les tranches s'alignent}, permettant au danger de traverser toutes les défenses.\\[0.3cm]
$\Rightarrow$ L'accident n'est \textbf{jamais} dû à une seule cause, mais à un \textbf{alignement} de facteurs (organisationnels + supervision + conditions + actes dangereux).
}}
\end{center}

\begin{memobox}
\textbf{80\% des accidents} en aviation sont liés aux \textbf{facteurs humains}. Le CRM a drastiquement réduit le taux d'accidents depuis les années 1980.
\end{memobox}

%=============================================================
\section{STRATÉGIE ENVIRONNEMENTALE --- AIR FRANCE ACT}
%=============================================================

\begin{tabular}{p{5cm}p{9.5cm}}
\toprule
\textbf{Programme} & \textbf{Air France ACT} \\
\midrule
Objectif 2030 (CO$_2$) & \textbf{$-$30\% par passager-km} vs 2019 (validé SBTi) \\
Objectif 2050 & Net zéro émissions \\
Incorporation SAF 2030 & \textbf{10\%} au niveau mondial \\
Flotte nouvelle génération & \textbf{70\%} d'ici 2030 (A220, A350, A320neo) \\
Investissement flotte & $>$1 Md\euro/an \\
Éco-pilotage & Montée continue, descente continue, roulage un moteur \\
Compensation carbone & Programme volontaire pour les passagers \\
\bottomrule
\end{tabular}

%=============================================================
\section{SYSTÈMES AVION}
%=============================================================

\subsection{Circuits Principaux d'un Avion de Ligne}

\begin{tabular}{p{3.5cm}p{4cm}p{7cm}}
\toprule
\textbf{Système} & \textbf{Source} & \textbf{Fonction} \\
\midrule
\textbf{Hydraulique} & 3 circuits (Vert/Bleu/Jaune sur Airbus) & Commandes de vol, train, freins, reverses \\
\textbf{Électrique} & Générateurs moteurs + APU + RAT (éolienne secours) & Instruments, éclairage, calculateurs, galley \\
\textbf{Pneumatique} & Prélèvement d'air moteurs (bleed air) & \textbf{Pressurisation cabine}, antigivrage, démarrage moteurs \\
\textbf{Carburant} & Réservoirs ailes + central & Alimentation moteurs, centrage (transfert) \\
\textbf{Oxygène} & Bouteilles (équipage) + générateurs chimiques (pax) & Urgence dépressurisation (masques tombent) \\
\bottomrule
\end{tabular}

\subsection{Pressurisation}
\begin{itemize}[nosep]
\item Cabine pressurisée à $\sim$\textbf{6\,000--8\,000 ft} d'altitude cabine (même si l'avion est à FL390)
\item Altitude cabine max certifiée~: 8\,000 ft (2\,438 m)
\item Au-dessus de \textbf{10\,000 ft} cabine (~3\,000 m)~: risque d'hypoxie $\Rightarrow$ masques à oxygène
\item \textbf{En cas de dépressurisation}~: descente d'urgence vers FL100 (10\,000 ft)
\end{itemize}

\subsection{APU (Auxiliary Power Unit)}
\begin{itemize}[nosep]
\item Petit turboréacteur auxiliaire dans la queue de l'avion
\item Fournit~: \textbf{électricité + air comprimé} au sol (sans moteurs principaux)
\item Permet le démarrage des moteurs principaux
\item Peut être utilisé en vol comme source de secours
\end{itemize}

%=============================================================
\section{BALISAGE ET MARQUAGE DES PISTES}
%=============================================================

\subsection{Numérotation des Pistes}
\begin{itemize}[nosep]
\item Le numéro = \textbf{cap magnétique arrondi à la dizaine} divisé par 10
\item Exemple~: piste orientée cap 086° $\Rightarrow$ piste \textbf{09}. L'autre sens~: 086°+180° = 266° $\Rightarrow$ piste \textbf{27}
\item Pistes parallèles~: suffixe L (Left), R (Right), C (Center)
\item CDG~: 08L/26R et 09R/27L (les 2 grandes pistes de 4\,200 m)
\end{itemize}

\subsection{Feux de Piste}

\begin{tabular}{p{4cm}p{10.5cm}}
\toprule
Feux de bord de piste & \textbf{Blancs} (jaunes sur les derniers 600m) \\
Feux de seuil & \textbf{Verts} (début de piste) \\
Feux de fin de piste & \textbf{Rouges} \\
Feux d'axe de piste & Blancs, puis alternés rouge/blanc, puis rouges \\
PAPI (indicateur de pente) & 4 feux~: 2 blancs + 2 rouges = sur le plan (3°) \\
\bottomrule
\end{tabular}

\begin{memobox}
\textbf{PAPI}~: 4 rouges = trop bas (danger !). 4 blancs = trop haut. 2 rouges + 2 blancs = correct (plan de 3°).
\end{memobox}

\subsection{Taxiways}
\begin{itemize}[nosep]
\item Ligne de roulage~: \textbf{jaune} (axe central)
\item Point d'arrêt avant piste~: \textbf{2 lignes jaunes continues + 2 tiretées} (NE JAMAIS franchir sans clairance ATC)
\end{itemize}

%=============================================================
\section{PHRASÉOLOGIE RADIO --- EXEMPLES}
%=============================================================

\subsection{Exemples de Communications}

\begin{tabular}{p{3cm}p{11.5cm}}
\toprule
\textbf{Phase} & \textbf{Exemple} \\
\midrule
Mise en route & \textit{``Air France 1234, request startup, information Bravo''} \\
Roulage & \textit{``Air France 1234, taxi to holding point runway 27R via Alpha, Bravo''} \\
Décollage & \textit{``Air France 1234, runway 27R, cleared for takeoff, wind 280/12''} \\
Montée & \textit{``Air France 1234, climb FL350''} \\
En route & \textit{``Air France 1234, maintain FL370, contact Paris Control 132.350''} \\
Descente & \textit{``Air France 1234, descend FL100, QNH 1018''} \\
Approche & \textit{``Air France 1234, cleared ILS approach runway 27R''} \\
Atterrissage & \textit{``Air France 1234, runway 27R, cleared to land, wind 260/08''} \\
\bottomrule
\end{tabular}

\subsection{Termes Importants}

\begin{tabular}{p{3cm}p{11.5cm}}
\toprule
\textbf{ROGER} & J'ai bien reçu votre message \\
\textbf{WILCO} & Will comply --- je vais exécuter \\
\textbf{AFFIRM} & Oui (ne pas dire ``yes'') \\
\textbf{NEGATIVE} & Non \\
\textbf{SAY AGAIN} & Répétez (ne pas dire ``repeat'', réservé à l'artillerie) \\
\textbf{MAYDAY} & Détresse (urgence vitale), répété 3 fois \\
\textbf{PAN PAN} & Urgence (non vitale), répété 3 fois \\
\textbf{SQUAWK} & Afficher un code transpondeur \\
\bottomrule
\end{tabular}

\begin{warnbox}
\textbf{MAYDAY} (3$\times$) = détresse vitale (feu, dépressurisation, panne moteur).\\
\textbf{PAN PAN} (3$\times$) = urgence non vitale (panne instrument, passager malade).\\
\textbf{Ne JAMAIS dire ``repeat''} en communication ATC~! On dit ``SAY AGAIN''.
\end{warnbox}

%=============================================================
\section{PERFORMANCES --- COMPLÉMENTS}
%=============================================================

\subsection{Effet du Vent sur les Performances}

\begin{tabular}{p{3cm}p{5.5cm}p{5.5cm}}
\toprule
\textbf{Type de vent} & \textbf{Au décollage} & \textbf{À l'atterrissage} \\
\midrule
\textbf{Vent de face} & $\downarrow$ distance de roulement (favorable) & $\downarrow$ distance d'atterrissage (favorable) \\
\rowcolor{warnred}
\textbf{Vent arrière} & $\uparrow$ distance de roulement (défavorable !) & $\uparrow$ distance d'atterrissage (défavorable !) \\
\textbf{Vent de travers} & Limitation max selon avion & Technique de décrabage \\
\bottomrule
\end{tabular}

\subsection{Effet de la Température et de l'Altitude}

\begin{tabular}{p{4cm}p{10.5cm}}
\toprule
\textbf{Température élevée} & Air moins dense $\Rightarrow$ moins de portance $\Rightarrow$ $\uparrow$ distance décollage, $\downarrow$ performances moteur \\
\textbf{Altitude terrain élevée} & Air moins dense (même effet). Aéroport en altitude = performances dégradées \\
\textbf{Altitude-densité} & Altitude fictive tenant compte de la vraie densité de l'air \\
\textbf{Piste mouillée} & $\uparrow$ distance de freinage. Risque d'aquaplaning \\
\bottomrule
\end{tabular}

\subsection{Distance Franchissable vs Autonomie}
\begin{itemize}[nosep]
\item \textbf{Distance franchissable}~: distance max en km (dépend de la vitesse, du vent, de la charge)
\item \textbf{Autonomie} (endurance)~: temps max en vol avec le carburant disponible
\item Le carburant embarqué inclut~: \textbf{trip fuel + réserves + déroutement + attente + taxi}
\item Réserve finale obligatoire~: \textbf{30 min} de vol (avion à réaction)
\end{itemize}

%=============================================================
\section*{QCM SUPPLÉMENTAIRES (Q101--120)}
\addcontentsline{toc}{section}{QCM SUPPLÉMENTAIRES (Q101--120)}
%=============================================================

\begin{enumerate}[nosep, start=101]
\item ETOPS-180 couvre quel \% du globe ? \\ a) 80\% \quad b) 90\% \quad c) 95\% \quad d) 100\%

\item Premier avion certifié ETOPS-370 ? \\ a) B777 \quad b) A350 \quad c) B787 \quad d) A330

\item Densité du Jet A-1 ? \\ a) 0.5 kg/L \quad b) 0.7 \quad c) $\sim$0.8 kg/L \quad d) 1.0

\item Point de congélation Jet A-1 ? \\ a) 0°C \quad b) $-$20°C \quad c) $-$38°C \quad d) $-$47°C

\item Le CRM a été développé après quel accident ? \\ a) Concorde \quad b) Tenerife 1977 \quad c) AF447 \quad d) 11 septembre

\item Modèle de Reason = ? \\ a) Dominos \quad b) Fromage suisse \quad c) Pyramide \quad d) Arbre

\item \% des accidents liés aux facteurs humains ? \\ a) 20\% \quad b) 50\% \quad c) $\sim$80\% \quad d) 100\%

\item Altitude cabine typique en croisière ? \\ a) 2\,000 ft \quad b) 6\,000--8\,000 ft \quad c) 15\,000 ft \quad d) 35\,000 ft

\item APU se trouve ? \\ a) Sous l'aile \quad b) Dans le cockpit \quad c) Dans la queue \quad d) Sous le fuselage

\item Numéro piste cap 270° ? \\ a) 09 \quad b) 27 \quad c) 18 \quad d) 36

\item PAPI 4 rouges = ? \\ a) Trop haut \quad b) Trop bas \quad c) Correct \quad d) Piste fermée

\item ``WILCO'' signifie ? \\ a) Roger \quad b) Will comply \quad c) Négatif \quad d) Répétez

\item MAYDAY = ? \\ a) Urgence \quad b) Détresse (vitale) \quad c) Panne radio \quad d) Information

\item SAF = ? \\ a) Safety Aviation Fuel \quad b) Sustainable Aviation Fuel \quad c) Standard Air Fuel \quad d) Secure Air Fuel

\item Objectif CO$_2$ AF 2030 ? \\ a) $-$10\% \quad b) $-$20\% \quad c) $-$30\% par pax-km \quad d) $-$50\%

\item Vent le plus favorable au décollage ? \\ a) Vent de face \quad b) Arrière \quad c) Travers \quad d) Aucun

\item Réserve finale min (jet) ? \\ a) 15 min \quad b) 30 min \quad c) 45 min \quad d) 60 min

\item Feux de seuil de piste = couleur ? \\ a) Rouge \quad b) Vert \quad c) Blanc \quad d) Jaune

\item Dépressurisation $\to$ descente vers ? \\ a) Sol \quad b) FL100 \quad c) FL200 \quad d) FL350

\item ``Say again'' remplace ? \\ a) Roger \quad b) Repeat \quad c) Affirm \quad d) Wilco
\end{enumerate}

\subsection*{Correction Q101--120}

\begin{center}
\begin{tabular}{|c|c||c|c||c|c||c|c|}
\hline
\textbf{Q} & \textbf{R} & \textbf{Q} & \textbf{R} & \textbf{Q} & \textbf{R} & \textbf{Q} & \textbf{R} \\
\hline
101 & c & 106 & b & 111 & b & 116 & a \\
102 & b & 107 & c & 112 & b & 117 & b \\
103 & c & 108 & b & 113 & b & 118 & b \\
104 & d & 109 & c & 114 & b & 119 & b \\
105 & b & 110 & b & 115 & c & 120 & b \\
\hline
\end{tabular}
\end{center}
